\rtmtight
\begin{tblr}{width=\textwidth,colspec={|Q[c,m,2.1cm]|X[l,m]|},hlines,vlines,hspan=minimal,rowsep=1.5pt,colsep=3pt,row{1}={bg=headgray,font=\bfseries}}
\SetCell{c,m}\hfil\logo\hfil & \textbf{POLITEKNIK NEGERI MALANG}\par\textbf{JURUSAN TEKNOLOGI INFORMASI}\par\textbf{PROGRAM STUDI: D4 TEKNIK INFORMATIKA} \\
\SetCell[c=2]{l} \textbf{RENCANA TUGAS MAHASISWA (RTM) DAN RUBRIK PENILAIAN} & \\
Nama Mata Kuliah & Proyek Inovasi \\
Kode Mata Kuliah & RTI256204 \quad \textbf{sks / Jam perminggu:} 6 SKS/12 jam \quad \textbf{Semester:} 6 \\
Dosen Pengampu & Tim Pengajar Program Studi D4 TI \\
\end{tblr}
\assessmentblock{Proposal Proyek Inovasi TI}{Case Method}{SCPMK0801-05301}{Mahasiswa menyusun proposal proyek inovasi TI yang mencakup analisis masalah industri, solusi teknologi, WBS, jadwal, anggaran, dan rencana manajemen risiko.}{Mengidentifikasi masalah nyata industri, merumuskan solusi teknologi, menyusun WBS dan timeline, menganalisis risiko, dan mempresentasikan proposal kepada dosen penguji.}{Dokumen proposal proyek, WBS, Gantt chart, matriks risiko, dan/atau bahan presentasi.}{Ketajaman analisis masalah dan relevansi solusi & 6 & Masalah industri dirumuskan dari data primer atau sekunder yang valid; solusi teknologi dipilih berdasarkan analisis kebutuhan yang sistematis dan berorientasi dampak. & Masalah dirumuskan dengan cukup jelas dan solusi relevan, tetapi analisis kebutuhan atau justifikasi pilihan teknologi masih dangkal. & Masalah tidak terdefinisi jelas, solusi tidak terkait kebutuhan industri, atau tidak ada justifikasi pilihan teknologi. \\ \\
Kelengkapan dan kualitas WBS serta jadwal & 5 & WBS mencakup seluruh deliverable, jadwal realistis dengan milestones jelas, dependensi antar-task terdefinisi, dan alokasi sumber daya sesuai kapasitas tim. & WBS dan jadwal cukup lengkap tetapi beberapa deliverable atau dependensi masih belum jelas. & WBS tidak mencakup semua deliverable, jadwal tidak realistis, atau tidak ada alokasi sumber daya. \\ \\
Kualitas manajemen risiko & 4 & Risiko teknis, tim, dan jadwal diidentifikasi secara komprehensif dengan probability, impact, dan strategi mitigasi yang operasional. & Risiko utama diidentifikasi tetapi probability, impact, atau mitigasi belum lengkap untuk semua risiko. & Risiko tidak diidentifikasi secara sistematis atau strategi mitigasi tidak operasional. \\}

\assessmentblock{Sprint Eksekusi Proyek Inovasi}{Project Based Learning}{SCPMK0803-05302, SCPMK0607-05304}{Mahasiswa mengeksekusi proyek inovasi melalui empat sprint agile, mengintegrasikan teknologi mutakhir, melaporkan kemajuan, mengelola perubahan, dan menghasilkan deliverable bertahap.}{Melaksanakan sprint planning, mengembangkan fitur sesuai backlog, melakukan sprint review dan retrospective, mengelola perubahan kebutuhan, dan mendokumentasikan setiap sprint.}{Backlog produk, sprint board, kode sumber terversionisasi, laporan sprint, dan demo artefak per sprint.}{Kualitas implementasi teknis dan inovasi solusi & 18 & Teknologi mutakhir diintegrasikan secara tepat dan inovatif; fitur utama berjalan, scalable, dan didukung oleh arsitektur yang terencana dengan baik. & Implementasi teknis berjalan untuk fitur utama tetapi integrasi teknologi atau aspek skalabilitas masih belum optimal. & Implementasi tidak menunjukkan inovasi teknologi, fitur utama belum berjalan, atau arsitektur tidak terencana. \\ \\
Kepatuhan sprint dan manajemen tim & 13 & Sprint dilaksanakan sesuai jadwal; backlog diperbarui setiap sprint; retrospective menghasilkan perbaikan nyata; perubahan dikelola dengan change log yang terdokumentasi. & Sprint sebagian besar sesuai jadwal, retrospective dilakukan, tetapi manajemen perubahan atau dokumentasi masih belum konsisten. & Sprint tidak sesuai jadwal, tidak ada retrospective, atau perubahan tidak dikelola secara sistematis. \\ \\
Kualitas dokumentasi dan deliverable sprint & 9 & Setiap sprint memiliki laporan kemajuan, demo artefak, dan kode terversionisasi dengan commit message yang informatif serta changelog yang lengkap. & Dokumentasi sprint ada tetapi beberapa laporan, demo, atau changelog masih tidak lengkap. & Dokumentasi sprint sangat minim atau artefak sprint tidak dapat diverifikasi. \\}

\assessmentblock{UTS Milestone Review}{UTS evaluasi tengah proyek}{SCPMK0803-05302}{Evaluasi tengah semester untuk mengukur kemajuan proyek inovasi, kualitas eksekusi sprint, dan kemampuan mahasiswa mengelola serta mempresentasikan status proyek.}{Mempresentasikan status proyek, menunjukkan demo MVP, menjelaskan tantangan yang dihadapi dan solusi yang diterapkan, serta rencana sprint berikutnya.}{Presentasi milestone, demo MVP/prototipe berjalan, laporan status proyek, dan backlog terbarui.}{Kemajuan proyek terhadap rencana & 6 & Kemajuan proyek sesuai atau melampaui rencana; deliverable yang dijanjikan tercapai; penyimpangan dijelaskan dengan analisis penyebab dan rencana koreksi yang konkret. & Kemajuan cukup sesuai rencana tetapi beberapa deliverable terlambat atau penyimpangan belum dianalisis dengan baik. & Kemajuan proyek jauh di bawah rencana tanpa analisis penyebab yang jelas atau rencana koreksi yang realistis. \\ \\
Kualitas demo dan presentasi milestone & 5 & Demo MVP berjalan dengan fitur utama yang dapat dioperasikan; presentasi menjelaskan arsitektur, keputusan teknis, dan rencana ke depan secara meyakinkan. & Demo berjalan untuk sebagian fitur; presentasi cukup jelas tetapi beberapa keputusan teknis atau rencana ke depan belum terjelaskan. & Demo tidak berjalan atau tidak representatif; presentasi tidak mampu menjelaskan status dan rencana proyek. \\ \\
Kemampuan mengelola risiko aktual & 4 & Risiko yang terjadi selama sprint diidentifikasi, dianalisis dampaknya, dan ditangani dengan tindakan konkret yang terdokumentasi dalam laporan status. & Risiko aktual diidentifikasi tetapi analisis dampak atau tindakan penanganan masih belum sistematis. & Risiko aktual tidak diidentifikasi atau tidak ada tindakan penanganan yang dapat diverifikasi. \\}

\assessmentblock{Laporan Akhir Proyek Inovasi}{Case Method}{SCPMK0804-05303}{Mahasiswa menyusun laporan akhir proyek inovasi yang mendokumentasikan seluruh proses pengembangan, hasil, evaluasi teknis dan bisnis, serta rekomendasi pengembangan lanjutan.}{Menyusun laporan komprehensif mencakup latar belakang, metodologi, hasil pengembangan, evaluasi, dan rekomendasi; merujuk pada semua artefak yang dihasilkan selama proyek.}{Laporan akhir tertulis, dokumentasi teknis sistem, hasil UAT, dan rekomendasi pengembangan.}{Kelengkapan dan sistematika laporan & 4 & Laporan mencakup semua komponen (latar belakang, metodologi, hasil, evaluasi, rekomendasi) dengan sistematika yang logis dan referensi artefak yang konsisten. & Laporan mencakup sebagian besar komponen tetapi beberapa bagian masih dangkal atau sistematikanya kurang konsisten. & Laporan tidak lengkap, tidak sistematis, atau tidak mencerminkan proses pengembangan yang sebenarnya. \\ \\
Kualitas evaluasi dan refleksi teknis & 4 & Evaluasi teknis mendalam mencakup analisis performa, scalability, keamanan, dan pembelajaran dari proses pengembangan; rekomendasi konkret dan dapat ditindaklanjuti. & Evaluasi teknis ada tetapi analisis kurang mendalam atau rekomendasi masih bersifat umum. & Evaluasi teknis tidak ada atau hanya berisi deskripsi tanpa analisis kritis. \\ \\
Kualitas dokumentasi teknis sistem & 2 & Dokumentasi teknis mencakup arsitektur, API, skema database, panduan deployment, dan cara penggunaan yang cukup untuk diserahterimakan kepada pihak industri. & Dokumentasi teknis ada tetapi tidak cukup lengkap untuk diserahterimakan tanpa penjelasan tambahan. & Dokumentasi teknis sangat minim atau tidak dapat digunakan sebagai acuan pengembangan lanjutan. \\}

\assessmentblock{UAS Defense Proyek Inovasi dan Demo}{UAS presentasi dan demo}{SCPMK0804-05303, SCPMK0607-05304}{Mahasiswa mempresentasikan dan mendemonstrasikan hasil proyek inovasi secara utuh kepada panel penguji, mempertahankan keputusan teknis dan bisnis, serta menunjukkan kesiapan produk untuk diimplementasikan.}{Mempresentasikan proyek inovasi secara komprehensif, mendemonstrasikan sistem berjalan, menjawab pertanyaan teknis dan bisnis dari penguji, dan menunjukkan nilai inovasi produk.}{Presentasi defense, demo sistem berjalan secara langsung, laporan akhir, dan artefak proyek lengkap.}{Kualitas demo dan kesiapan produk & 8 & Sistem berjalan sempurna dengan semua fitur utama, skalabel, dan siap diimplementasikan; demo menunjukkan alur penggunaan nyata sesuai kebutuhan industri yang diidentifikasi. & Sistem berjalan untuk fitur utama tetapi beberapa fitur pendukung belum selesai atau kesiapan produk masih memerlukan pengembangan lanjutan. & Sistem tidak dapat didemonstrasikan secara penuh atau produk tidak siap untuk diimplementasikan. \\ \\
Defense teknis dan inovasi solusi & 7 & Keputusan arsitektur dan teknologi dipertahankan dengan argumen teknis yang kuat; inovasi solusi dijelaskan secara meyakinkan termasuk keunggulan dibanding pendekatan konvensional. & Defense teknis cukup meyakinkan tetapi beberapa keputusan atau aspek inovasi belum dapat dipertahankan dengan argumen yang kuat. & Defense teknis tidak dapat mempertahankan keputusan utama atau nilai inovasi tidak dapat dijelaskan. \\ \\
Nilai bisnis dan dampak solusi & 5 & Nilai bisnis solusi dijelaskan dengan data atau estimasi yang terukur; dampak terhadap masalah industri yang diidentifikasi terbukti dan dapat diukur. & Nilai bisnis dijelaskan tetapi estimasi atau bukti dampak masih bersifat kualitatif dan belum terukur. & Nilai bisnis tidak dijelaskan atau tidak ada keterkaitan dengan masalah industri yang diidentifikasi. \\}

\begin{tblr}{width=\textwidth,colspec={|Q[l,m,3.2cm]|X[l,m]|},hlines,vlines,hspan=minimal,row{1-3}={bg=headgray},rowsep=1.5pt,colsep=3pt}
Jadwal Pelaksanaan: & Minggu 1--17 sesuai RPS. Setiap evaluasi memiliki RTM dan rubrik multi-kriteria. \\
Lain-Lain Yang Diperlukan: & LMS, repositori proyek (Git), project management tools, perangkat lunak sesuai mata kuliah, dan perangkat presentasi. \\
Pustaka: & Mengacu pustaka utama dan pendukung pada bagian RPS. \\
\end{tblr}
